% greeninfer_ecoquery.tex
% Compile: pdflatex greeninfer_ecoquery.tex  (run twice for references)

\documentclass[11pt,a4paper]{article}

\usepackage[T1]{fontenc}
\usepackage[utf8]{inputenc}
\usepackage{lmodern}
\usepackage[margin=2.5cm]{geometry}
\usepackage{hyperref}
\usepackage{booktabs}
\usepackage{tabularx}
\usepackage{array}
\usepackage{amsmath}
\usepackage{xcolor}
\usepackage{listings}
\usepackage{microtype}
\usepackage{authblk}
\usepackage{abstract}
\usepackage{parskip}
\usepackage{setspace}
\usepackage[font=it,labelfont=it,position=below]{caption}

\hypersetup{
  colorlinks = true,
  linkcolor  = blue!70!black,
  citecolor  = blue!70!black,
  urlcolor   = blue!70!black,
}

\lstset{
  basicstyle      = \ttfamily\small,
  breaklines      = true,
  frame           = single,
  backgroundcolor = \color{gray!8},
}

\title{\textbf{GreenInfer + EcoQuery: Estimating the CO\textsubscript{2} and Water
       Footprint of AI Inference Queries}}

\author{Gayani Upeksha}
\date{November 2026}

\begin{document}
\maketitle
\thispagestyle{empty}

\begin{abstract}
The rapid growth of AI services has raised growing concerns about their
environmental cost. Each inference query to a large language model consumes
measurable electricity and produces CO\textsubscript{2} emissions and water
consumption that vary considerably by location and local energy mix. Most
existing estimation tools either target training workloads, lack regional
detail, or ignore water footprint entirely. This paper introduces
\textbf{GreenInfer}, a Python library, and \textbf{EcoQuery}, an interactive
web dashboard, that together estimate the CO\textsubscript{2} emissions and
water footprint of a single AI inference query across 209 countries and
all 31 Chinese provinces. Country grid intensity draws on Ember data
spanning 2000--2025, enabling historical and cross-year comparison even
though AI inference at scale only became common from 2022 onward
(2025 data are partial, covering approximately 93 countries at time of writing);
provincial data covers 2024--2025. Provincial carbon
intensity is derived from National Bureau of Statistics (NBS) monthly
generation data disaggregated by source, using the IPCC AR5 lifecycle
emission factor for thermal generation. Water footprint follows the
Water Usage Effectiveness (WUE) methodology of Ren et al.\ (2025). Results show that the 2025 per-query CO2 emissions range from
26~gCO2 (Xizang) to 787~gCO2 (Beijing), while water consumption
is approximately 4.9~mL per query regardless of location. Both tools are open-source.
\end{abstract}

\vspace{1em}
\noindent\textbf{Keywords:} AI inference, carbon footprint, water footprint,
China electricity, carbon intensity, sustainability

\section{Introduction}

Large language models have become routine infrastructure. Since ChatGPT's
public launch in late 2022, AI inference---generating a response to a user
query---has grown into one of the most common computational workloads
worldwide. Unlike training, which happens once and has attracted substantial
research attention, inference runs billions of times per day and its
cumulative environmental cost is poorly understood.

A single AI inference query consumes roughly 0.001~kWh of
electricity~\cite{patterson2021}, around ten times the energy of a
traditional web search. The resulting CO\textsubscript{2} emissions depend
on the carbon intensity of the grid powering the data center, which varies
substantially by country, region, and year. A query served from a
coal-heavy data center in Shanxi produces far more emissions than the same
query served from a hydropower-dominated grid in Sichuan. Tools such as
CodeCarbon~\cite{codecarbon}, ML CO\textsubscript{2}
Impact~\cite{mlco2impact}, and Green Algorithms~\cite{greenalgorithms}
estimate emissions for training and batch workloads, but none offer
per-query estimates at the Chinese provincial level, and none integrate
water footprint.

Water is an overlooked dimension of AI's environmental impact. Data centers
use water for cooling, and electricity generation itself consumes water
upstream. Ren et al.\ (2025)~\cite{ren2025} developed a WUE-based framework
separating on-site and off-site water consumption, making per-query water
estimation tractable.

This paper contributes:
\begin{enumerate}
  \item \textbf{GreenInfer}: an open-source Python library estimating
        CO\textsubscript{2} and water footprint per AI inference query for
        any country or Chinese province.
  \item \textbf{EcoQuery}: an interactive dashboard covering 209 countries
        (2000--2025) and 31 Chinese provinces (2024--2025).
  \item A provincial carbon intensity dataset derived from NBS monthly
        generation data, the first calibrated to the AI inference era.
  \item Integration of the Ren et al.\ (2025) WUE methodology at regional
        granularity.
\end{enumerate}

\section{Related Work}

\paragraph{AI energy and carbon tools.}
CodeCarbon~\cite{codecarbon} measures real-time energy consumption of Python
code and maps it to CO\textsubscript{2} using regional grid intensity.
ML CO\textsubscript{2} Impact~\cite{mlco2impact} offers a web calculator for
training emissions based on hardware and cloud provider.
Green Algorithms~\cite{greenalgorithms} extends this to HPC workloads. All
three share a fundamental limitation: they target training and batch jobs,
not individual inference queries, and none resolve China at the provincial
level.

\paragraph{China electricity carbon intensity.}
The CEADs project~\cite{ceads} provides provincial CO\textsubscript{2}
inventories for China from 2000 to 2019, covering energy, industry, and
agriculture. The 2019 cutoff predates large-scale LLM deployment, and
CEADs does not disaggregate electricity generation by source at the
provincial level in a way suitable for intensity calculations. The NBS
publishes monthly provincial generation broken down by source (thermal,
hydro, nuclear, wind, solar), providing the most current and granular data
available for this purpose.

\paragraph{Water footprint of AI.}
Ren et al.\ (2025)~\cite{ren2025} introduced a systematic framework for
AI inference water footprint, distinguishing on-site cooling water from
off-site water consumed during electricity generation. Their WUE
methodology provides the first peer-reviewed basis for per-query water
estimation. Li et al.\ (2023)~\cite{li2023} estimated that training GPT-3
consumed approximately 700,000 liters of water, drawing wider attention to
AI's water impact.

\paragraph{Gap addressed.}
No existing tool combines per-query CO\textsubscript{2} estimation,
Chinese provincial granularity using current NBS data, and water footprint
via the Ren et al.\ (2025) WUE framework. GreenInfer and EcoQuery fill
this gap.

\section{Methodology}

\subsection{Energy per Query}

Following Patterson et al.\ (2021)~\cite{patterson2021}, we use a
per-query energy constant of $E = 0.001$~kWh. This reflects typical
transformer inference on modern GPU hardware and is consistent with widely
cited values in the literature.

\subsection{CO\textsubscript{2} Emissions}

Per-query CO\textsubscript{2} is:
\begin{equation}
  \text{CO}_2 \; (\text{g}) = E \times I
\end{equation}
where $E$ is energy per query (kWh) and $I$ is grid carbon intensity
(gCO\textsubscript{2}/kWh) at the data center location.

\paragraph{World countries.}
For 209 countries we use Ember Yearly Electricity Data~\cite{ember}, which
provides emissions intensity values per country from 2000 to 2025.
Although large-scale AI inference only became common from late 2022, we
include the full 2000--2025 range for two reasons. First, the grid carbon
intensity of a country in any given year is a property of its electricity
system, independent of whether AI workloads existed at that time; a
researcher or policymaker may legitimately ask what a hypothetical inference
query \emph{would} have cost had it been run in 2010. Second, historical
coverage enables trend analysis---showing how rapidly a country's grid has
decarbonised---which contextualises current per-query footprints. Users who
wish to restrict analysis to the AI era can simply query years from 2022
onward.

\paragraph{Chinese provinces.}
For all 31 provinces we derive intensity from NBS monthly generation data
disaggregated by source. The formula is:
\begin{equation}
  I_{\text{province}} =
  \frac{\text{Thermal}_{\text{annual}} \times 820}%
       {(\text{Thermal} + \text{Hydro} + \text{Nuclear} + \text{Wind} +
         \text{Solar})_{\text{annual}}}
\end{equation}

The factor 820~gCO\textsubscript{2}/kWh is the IPCC AR5 lifecycle emission
factor for coal~\cite{ipcc2014}. China's thermal generation is approximately
95\% coal, making this assumption reasonable at the provincial level. Zero
emissions are assigned to hydro, nuclear, wind, and solar.

Data cover 2024 and 2025. The NBS combines January and February figures and
releases them in March, so these two months are absent from individual
monthly files. Annual totals are estimated from 10 available months scaled
by 12/10; this assumption is discussed as a limitation in
Section~\ref{sec:limitations}.

\subsection{Water Footprint}

We follow the WUE framework of Ren et al.\ (2025)~\cite{ren2025}:
\begin{equation}
  \text{water} \; (\text{mL}) =
  \left(\frac{E}{\text{PUE}} \times \text{WUE}_{\text{onsite}}
   + E \times \text{WUE}_{\text{offsite}}\right) \times 1000
\end{equation}

On-site WUE captures data center cooling; off-site WUE captures water
consumed during electricity generation. Parameter values:

\begin{table}[h]
\centering
\renewcommand{\arraystretch}{1.2}
\begin{tabular}{lll}
\toprule
Parameter & Value & Source \\
\midrule
$E$ & 0.001~kWh & Patterson et al.\ (2021) \\
PUE & 1.4 & Industry average \\
WUE\textsubscript{onsite}  & 0.25~L/kWh & Ren et al.\ (2025) \\
WUE\textsubscript{offsite} & 4.73~L/kWh & Ren et al.\ (2025) \\
\bottomrule
\end{tabular}
\vspace{4pt}
\caption{Water footprint model parameters.}
\end{table}

This gives approximately \textbf{4.91~mL per query}, uniform across all
regions since the parameters reflect data center and grid averages rather
than region-specific values.

\subsection{Data Sources}

\begin{table}[h]
\centering
\renewcommand{\arraystretch}{1.2}
\begin{tabular}{lll}
\toprule
Data & Source & Coverage \\
\midrule
Country electricity intensity   & Ember~\cite{ember}      & 209 countries, 2000--2025 \\
China provincial generation     & NBS~\cite{nbs}          & 31 provinces, 2024--2025 \\
Thermal emission factor         & IPCC AR5~\cite{ipcc2014}& 820~gCO\textsubscript{2}/kWh \\
Water footprint parameters      & Ren et al.~\cite{ren2025} & WUE onsite + offsite \\
\bottomrule
\end{tabular}
\vspace{4pt}
\caption{Data sources.}
\end{table}

\section{Results}

\subsection{Chinese Provincial CO\textsubscript{2} Intensity}

Table~\ref{tab:provinces} shows estimated intensity for all 31 provinces in
2024 and 2025. The range is wide, reflecting the stark variation in China's
regional energy mix. ($\Delta$ = change from 2024 to 2025.)

\begin{table}[htbp]
\centering
\renewcommand{\arraystretch}{1.2}
\label{tab:provinces}
\small
\begin{tabular}{lrrr}
\toprule
Province & 2024 & 2025 & $\Delta$ \\
\midrule
Beijing        & 795 & 787 & $-$8  \\
Shanghai       & 794 & 784 & $-$10 \\
Tianjin        & 748 & 729 & $-$19 \\
Anhui          & 738 & 719 & $-$19 \\
Shaanxi        & 713 & 701 & $-$12 \\
Jiangxi        & 665 & 670 & $+$5  \\
Shanxi         & 687 & 663 & $-$24 \\
Jiangsu        & 666 & 653 & $-$13 \\
Shandong       & 686 & 651 & $-$35 \\
Henan          & 676 & 650 & $-$26 \\
Zhejiang       & 606 & 611 & $+$5  \\
Chongqing      & 646 & 610 & $-$36 \\
Ningxia        & 618 & 593 & $-$25 \\
Guangdong      & 583 & 589 & $+$6  \\
Inner Mongolia & 624 & 579 & $-$45 \\
Heilongjiang   & 585 & 575 & $-$10 \\
Hebei          & 611 & 567 & $-$44 \\
Xinjiang       & 595 & 561 & $-$34 \\
Jilin          & 488 & 489 & $+$1  \\
Hunan          & 471 & 483 & $+$12 \\
Hainan         & 556 & 479 & $-$77 \\
Fujian         & 478 & 476 & $-$2  \\
Guizhou        & 538 & 463 & $-$75 \\
Liaoning       & 463 & 455 & $-$8  \\
Gansu          & 412 & 429 & $+$17 \\
Hubei          & 411 & 370 & $-$41 \\
Guangxi        & 393 & 347 & $-$46 \\
Sichuan        & 147 & 145 & $-$2  \\
Yunnan         & 107 &  97 & $-$10 \\
Qinghai        &  99 & 103 & $+$4  \\
Xizang (Tibet) &  23 &  26 & $+$3  \\
\bottomrule
\end{tabular}
\vspace{4pt}
\caption{CO\textsubscript{2} intensity (gCO\textsubscript{2}/kWh) for all
31 Chinese provinces, 2024 and 2025. Derived from NBS monthly generation
data using the IPCC AR5 coal factor of 820~gCO\textsubscript{2}/kWh.
Sorted by 2025 intensity, descending.}
\end{table}

The highest intensities are in eastern industrial provinces where thermal
generation dominates: Beijing (787), Shanghai (784), and Tianjin (729) in
2025. The lowest are in the hydropower-rich southwest: Xizang, Qinghai, and
Sichuan all fall below 150~gCO\textsubscript{2}/kWh. The gap between
Beijing and Xizang---roughly 30$\times$---illustrates how much geographic
placement matters for the environmental cost of AI services.

\subsection{International Comparison}

Using Ember data for 2024, per-query CO\textsubscript{2} spans an even
wider range internationally. Iceland, Norway, and Paraguay approach
zero due to near-total renewable generation. South Africa, Poland, and
Kazakhstan are among the highest due to coal dependence. China's national
average (around 550~gCO\textsubscript{2}/kWh in 2024) sits in the
upper-middle range; the United States (around 370) and Germany (around 290)
are considerably lower as a result of recent renewables growth.

\subsection{Year-on-Year Trend}

Comparing 2024 to 2025 across all 31 provinces, the dominant pattern is
modest decline. Inner Mongolia ($-45$), Guangxi ($-46$), Hebei ($-44$),
Chongqing ($-36$), and Hainan ($-77$) recorded the largest reductions,
consistent with accelerating wind and solar capacity additions in those
regions. A handful of provinces show slight increases---Jiangxi (+5),
Zhejiang (+5), Guangdong (+6), Gansu (+17), Hunan (+12)---likely reflecting
inter-annual variation in hydro availability and load growth outpacing new
generation in those grids. Xizang remains the cleanest at
23--26~gCO\textsubscript{2}/kWh across both years, powered almost entirely
by run-of-river hydropower.

\subsection{Water Footprint}

Each inference query consumes approximately \textbf{4.91~mL of water}:
\begin{itemize}
  \item \textbf{0.18~mL} on-site (data center cooling):
        $0.001 / 1.4 \times 0.25 \times 1000$
  \item \textbf{4.73~mL} off-site (electricity generation):
        $0.001 \times 4.73 \times 1000$
\end{itemize}

Off-site water accounts for 96\% of the total. This means that efficiency
improvements at the data center level have limited impact on overall water
footprint; reductions require cleaner electricity grids or lower-water
generation technologies upstream.

\section{Discussion}

\subsection{Limitations}
\label{sec:limitations}

\paragraph{January-February scaling.}
The NBS releases January and February provincial generation figures combined,
delayed until March, so both months are absent from individual monthly
files. Annual totals for 2024 and 2025 are estimated from 10 months scaled
by 12/10. This approach assumes the missing months have a proportional
generation mix. In practice, January-February coincides with peak heating
demand in northern provinces, which may push the thermal share slightly
higher than the year-round average. Future work should substitute the
released combined figures once available.

\paragraph{Thermal emission factor.}
We apply a uniform 820~gCO\textsubscript{2}/kWh factor to all thermal
generation in all provinces. This is appropriate given that China's thermal
mix is around 95\% coal nationally. However, some provinces have meaningful
gas-fired capacity (emission factor around 490~gCO\textsubscript{2}/kWh),
and the provincial thermal fuel split is not publicly available in NBS data.
This may cause slight overestimation in provinces with significant gas.

\paragraph{Fixed energy per query.}
The 0.001~kWh figure is a literature average for large transformer models on
modern GPU hardware. Actual energy varies considerably with model size,
hardware generation, and batch configuration. GreenInfer accepts a custom
\texttt{energy\_kwh} parameter to mitigate this for users with hardware
measurements.

\paragraph{Water footprint uniformity.}
The WUE parameters are global averages. Regional values would improve
accuracy---water consumption per kWh of generation depends on technology
(hydro vs.\ gas cooling vs.\ coal) and local climate---but are not currently
available at the provincial level.

\subsection{Absolute vs.\ Intensity-Based Trends}

A key nuance in interpreting the national trend data is the distinction
between carbon intensity (gCO\textsubscript{2}/kWh) and total carbon
emissions. China's grid carbon intensity has declined steadily---from
approximately 670~gCO\textsubscript{2}/kWh in 2011 to around
491~gCO\textsubscript{2}/kWh in 2025---because renewable capacity (wind,
solar, hydro) is growing faster than total demand. However, thermal
generation has also increased in absolute terms over the same period, as
electricity consumption continues to rise with economic growth. The result
is that the \emph{per-query} carbon footprint of AI inference in China is
improving year by year, while \emph{total} grid CO\textsubscript{2}
emissions from electricity remain elevated. For AI specifically, this
means that falling per-query intensity does not automatically translate
to a falling absolute carbon footprint: if inference query volume grows
faster than the grid decarbonises---a plausible scenario given rapid AI
adoption---the aggregate environmental cost continues to rise. This
reinforces the importance of both grid transition and query-level
efficiency improvements.

\subsection{Future Work}

Several natural extensions follow from this work. Substituting released NBS
January-February figures for the 12/10 scaling estimate would improve
annual accuracy. Extending provincial data to 2026 as records become
complete would enable a three-year trend analysis. Developing
province-specific WUE parameters would substantially improve water
estimates. Integrating a live grid intensity API (such as Electricity Maps)
would enable real-time per-query estimation in place of annual averages.

\section{Conclusion}

This paper introduced GreenInfer and EcoQuery, tools for estimating the
CO\textsubscript{2} emissions and water footprint of a single AI inference
query. By combining Ember global electricity data with NBS provincial
generation data and the Ren et al.\ (2025) WUE methodology, the system
covers 209 countries (2000--2025) and 31 Chinese provinces (2024--2025),
and is the first to offer both dimensions in a single open-source package.

The key findings are straightforward: per-query carbon footprint varies by
more than 30$\times$ across Chinese provinces depending on the local grid,
and 96\% of the water footprint comes from electricity generation rather
than on-site cooling. Both results point to the same conclusion---where AI
infrastructure is located matters enormously, and water deserves the same
attention as carbon in AI sustainability assessments.

GreenInfer and EcoQuery are publicly available, enabling researchers,
developers, and policymakers to compare the environmental cost of AI
services across regions and over time.

\begin{thebibliography}{10}

\bibitem{patterson2021}
D.~Patterson et al.,
``Carbon considerations for large language model training,''
\emph{arXiv preprint arXiv:2104.10350}, 2021.

\bibitem{codecarbon}
V.~Lottick et al.,
``Energy usage reports: Environmental awareness as part of algorithmic
accountability,''
\emph{NeurIPS Workshop on Tackling Climate Change with ML}, 2019.

\bibitem{mlco2impact}
A.~Lacoste, A.~Luccioni, V.~Schmidt, and T.~Dandres,
``Quantifying the carbon emissions of machine learning,''
\emph{arXiv preprint arXiv:1910.09700}, 2019.

\bibitem{greenalgorithms}
L.~Lannelongue, J.~Grealey, and M.~Inouye,
``Green algorithms: Quantifying the carbon footprint of computation,''
\emph{Advanced Science}, vol.~8, no.~12, p.~2100707, 2021.

\bibitem{ren2025}
Y.~Ren et al.,
``Estimating the water footprint of AI model inference,''
\emph{arXiv preprint arXiv:2505.09598}, 2025.

\bibitem{ceads}
S.~Guan et al.,
``CEADs: A dataset of CO\textsubscript{2} emissions accounts for China,''
\emph{Scientific Data}, vol.~8, no.~1, p.~193, 2021.

\bibitem{li2023}
S.~Li et al.,
``Making AI less `thirsty': Uncovering and addressing the secret water
footprint of AI models,''
\emph{arXiv preprint arXiv:2304.03271}, 2023.

\bibitem{ember}
Ember,
``Yearly Electricity Data,''
\emph{Ember Climate}, 2025.
[Online]. Available: \url{https://ember-climate.org/data/data-tools/data-explorer/}

\bibitem{nbs}
National Bureau of Statistics of China,
``Monthly By Province: Output of Energy Products,''
\emph{NBS Data}, 2024--2025.
[Online]. Available: \url{https://data.stats.gov.cn}

\bibitem{ipcc2014}
O.~Edenhofer et al.\ (Eds.),
\emph{Climate Change 2014: Mitigation of Climate Change.
Contribution of Working Group~III to the Fifth Assessment Report of the IPCC},
Cambridge University Press, 2014.

\end{thebibliography}

\end{document}
